\subsubsection{Coordenadas centrales, orientación y registros regionales}
\label{el-centro:registros}

La estructura central relaciona posición, evaluación y transporte. En la
celda \((5,5)\), las dos hojas de APP dan
\begin{equation}
 \mathcal A(5,5)=\bigl((1,1),(7,2)\bigr),\qquad
 5+5=1+9\cdot1,\qquad5\cdot5=7+9\cdot2.
 \label{el-centro:app}
\end{equation}
Los dos unos de la hoja aditiva son, en su orden, residuo y cociente.
Su coincidencia numérica conserva dos funciones aritméticas distintas.
Sobre la fibra de suma \(10\), las dos vacancias orientadas verifican
\begin{equation}
 \mathcal A(8,2)=\mathcal A(2,8)=\bigl((1,1),(7,1)\bigr).
 \label{el-centro:vacancias}
\end{equation}
Se conservan la hoja aditiva completa y el residuo multiplicativo;
la coordenada de cociente del producto disminuye exactamente en
\(1\). El transporte retiene, además, cuál de las dos posiciones
conjugadas se ha alcanzado.

\begin{figure}[H]
\centering
\begin{tikzpicture}[x=.67cm,y=.67cm,font=\footnotesize,>=Latex]
 \draw[step=1,black!12,line width=.3pt] (1,1) grid (9,9);
 \draw[->,black!60] (.7,1)--(9.5,1) node[right] {$i$};
 \draw[->,black!60] (1,.7)--(1,9.5) node[above] {$j$};
 \foreach \k in {1,...,9}{
  \node[below=2pt] at (\k,1) {$\k$};
  \node[left=2pt] at (1,\k) {$\k$};
 }
 \draw[black!65] (1,9)--(9,1);
 \draw[<->,line width=.9pt] (2,8)--(8,2);
 \foreach \x/\y in {2/8,5/5,8/2}{
  \filldraw[fill=white,line width=.8pt] (\x,\y) circle (3pt);
 }
 \fill (5,5) circle (1.5pt);
 \node[above right=3pt,fill=white,inner sep=1pt] at (2,8) {$(2,8)$};
 \node[above right=3pt,fill=white,inner sep=1pt] at (5,5) {$(5,5)$};
 \node[above right=3pt,fill=white,inner sep=1pt] at (8,2) {$(8,2)$};
 \node[anchor=west,align=left] at (10.5,7.6)
   {$i+j=10$\\$R^+=1,\quad q^+=1$};
 \node[anchor=west,align=left] at (10.5,5.4)
   {centro: $ij=25$\\$R^\times=7,\quad q^\times=2$};
 \node[anchor=west,align=left] at (10.5,3.1)
   {vacancias: $ij=16$\\$R^\times=7,\quad q^\times=1$};
\end{tikzpicture}
\caption{Centro electrónico y vacancias conjugadas en la coordenatización APP.
La diagonal es la fibra aditiva \(\mathcal F_{10}\); el centro
es el punto fijo de \((i,j)\mapsto(10-i,10-j)\). Los dos
desplazamientos de módulo \(3\) conservan la suma y el residuo
producto y reducen en una unidad el cociente multiplicativo. Las
flechas representan las deformaciones sobre la fibra.}
\label{el-centro:figura}
\end{figure}


La simetría central caracteriza también su persistencia bajo cambios de
coordenatización. Si una transformación admisible \(f\) del atlas entrelaza la
inversión, \(f\rho=\rho f\), entonces
\[
 \rho f(5,5)=f\rho(5,5)=f(5,5).
\]
La unicidad del punto fijo obliga a \(f(5,5)=(5,5)\). Ésta es la
propiedad interna que hace estable la localización central: toda
transformación de ese tipo conserva la celda antes de evaluar su
carácter electrónico (sección~\ref{el-centro:registros}).

El bloque matricial central pertenece a la evaluación sobre las dos
posiciones adyacentes \(4,5\). Su descomposición, establecida en el
desarrollo local del núcleo, es
\[
 B_c=\begin{pmatrix}7&2\\2&7\end{pmatrix}=9P_++5P_-.
\]
Aquí \(9\) y \(5\) son autovalores de dos modos, mientras que
\((5,5)\) es una posición del atlas. Las concatenaciones de sus
entradas dan \(72+27=77+22=99\) y
\(77-22=55=54+1\). La primera igualdad conserva la suma total bajo
dos disposiciones; la segunda expresa el defecto entre las lecturas
diagonal y antidiagonal. Estas operaciones preceden a la composición
electrónica y conservan su procedencia matricial.

La orientación de partida es otro dato. En la representación del
retorno electrónico se fija \((h_0,v_0)=(1,1)\), es decir, avance
positivo en ambas coordenadas. La inversión simultánea cada \(27\)
pasos, leída en ventanas de \(9\), determina
\[
 \tau_e=(+1,+1,+1,-1,-1,-1,+1,+1,+1,-1,-1,-1).
\]
Los tramos \((1,1,1)\) pertenecen a esta sucesión de orientaciones:
cada tramo comprende tres ventanas del mismo sentido. La recurrencia
de la sección de lectores demuestra este registro en
\eqref{eq:el-tau}, y sus dos evaluaciones recuperan el triple
\((80,54,6)\) que interviene en
\eqref{eq:el-formula-completa}. La coordenada de posición, los
cocientes APP y la sucesión de signos quedan así conectados mediante
operaciones explícitas.

La firma \(555555\) pertenece a una lectura regional de otro
dominio. Para una emisión decimal \(U=(U_1,\ldots,U_6)\), escribimos
\(U_j=100d_{1j}+10d_{2j}+d_{3j}\), con tres cifras, incluidos los
ceros iniciales. El lector de firma multiplicativa es
\[
 \mathcal E_\times(U)
   =\bigl([d_{2j}-d_{1j}]_{10}\bigr)_{j=1}^{6}.
\]
En la región transversal
\(U_\pi^\perp=(501,614,498,169,272,272)\), las diferencias son
\((-5,-5,5,5,5,5)\). Por consiguiente,
\begin{equation}
 \mathcal E_\times(U_\pi^\perp)=(5,5,5,5,5,5),\qquad
 \Psi(U_\pi^\perp)=010211.
 \label{el-centro:firma-transversal}
\end{equation}
La primera aplicación conserva una firma multiplicativa del catálogo;
la segunda conserva la palabra ternaria común a las cinco regiones.
La figura \ref{pi-estrella:figura} muestra la posición transversal de
esta emisión y su multiplicidad \(432\).

\paragraph{Memoria necesaria para transportar la firma constante.}
La firma \(555555\) proporciona, además, un ejemplo verificable de
la información de trayectoria que exige TPK. Sea
\(\mathcal X_\times=D_9^2\times\{N,E,S,O\}\), con \(324\)
posiciones orientadas del canal producto. El avance \(P\) desplaza
una casilla en la dirección almacenada y la media vuelta \(H\)
invierte esa dirección. En esta representación, el lector
\(E_{\rm sig}\) suma por ventanas los exponentes discretos
\[
 \ell_9(1,2,4,8,7,5)=(0,1,2,3,4,5),\qquad
 \ell_9(3)=\ell_9(6)=\ell_9(9)=0.
\]
Durante cada una de las seis ventanas de nueve pasos, se evalúa
\(\ell_9(\rho_9(ij))\) antes de cada avance producto, seleccionado
por TRIT en los tiempos \(2,5,8\pmod9\). La suma de la ventana se
reduce módulo \(10\). Tras los pasos \(27\) y \(54\) se invierte
la dirección. Esta regla determina las seis componentes de
\(E_{\rm sig}\) desde la posición y orientación iniciales.

La enumeración de este dominio produce \(26\) firmas. La fibra de
\(555555\) contiene \(24\) posiciones orientadas; después de un
avance, sus imágenes se distribuyen en tres firmas, con ocho
representantes cada una:
\[
 555177,\qquad555717,\qquad555771.
\]
Tres testigos, en coordenadas de \(1\) a \(9\), permiten comprobar
la distinción directamente:
\[
\begin{array}{c|c|c}
x&E_{\rm sig}(x)&E_{\rm sig}(Px)\\\hline
(3,5,N)&555555&555177\\
(3,4,N)&555555&555717\\
(3,4,S)&555555&555771
\end{array}
\]
Por tanto, el valor común de la firma no determina su siguiente valor:
la clase de transporte dentro de esa fibra forma parte del dato
operativo. La partición estable mínima se calcula partiendo de la
equivalencia de firmas \(\sim_0\) y refinando mediante
\[
 x\sim_{n+1}y\quad\Longleftrightarrow\quad
 x\sim_n y,\quad Px\sim_n Py,\quad Hx\sim_n Hy.
\]
El procedimiento termina al ser finito el dominio. Toda congruencia
estable que refine \(\sim_0\) refina cada \(\sim_n\), por inducción;
su punto fijo es, por ello, la congruencia estable mínima requerida.
El censo da \(26\to28\to28\), con histograma
\(27\cdot8+108=324\). La única firma que se descompone es
\(555555\), en tres clases de ocho elementos. El certificado adjunto
reproduce la enumeración, las tres clases y la estabilidad por ambos
operadores (sección~\ref{el-centro:registros}).

La multiplicidad \(24\) pertenece aquí a la fibra del canal producto;
la multiplicidad \(432\) pertenece a la emisión regional completa
\(U_\pi^\perp\). La primera prueba por qué la firma necesita memoria
de transporte; la segunda sitúa esa firma en el catálogo conjunto de
las dos hojas. Ambas lecturas participan de la misma arquitectura,
con dominios y operaciones especificados.

Esta distinción de dominios permite reunir las apariciones de la unidad
y del centro con sus respectivos significados:
\begin{center}
\small
\begin{tabularx}{\textwidth}{@{}lX@{}}
\toprule
Registro & Objeto y función \\
\midrule
\((5,5)\) & Celda fija de la inversión central del atlas APP.\\
\((R^+,q^+)=(1,1)\) & Residuo y cociente de la evaluación aditiva central.\\
\((h_0,v_0)=(1,1)\) & Orientación inicial de los dos desplazamientos.\\
\((+1,+1,+1)\) & Tres ventanas consecutivas de \(\tau_e\) con orientación positiva.\\
\(c=111111\) & Dato de la quinta frontera nonádica, sección \ref{mono-recta:desarrollo}.\\
\(\mathcal E_\times=555555\) & Firma de la región transversal de \(\pi\), con seis posiciones.\\
\bottomrule
\end{tabularx}
\end{center}

La coordinación electrónica se obtiene componiendo estos niveles en su
orden de dependencia. APP determina el centro y las evaluaciones;
TRIT selecciona regímenes y orientaciones; TPK prolonga la trayectoria
y conserva sus registros. Las regiones previamente generadas aportan
los caracteres \(\pi,\varphi,e\); el registro dodecafásico aporta
\(\alpha\); los lectores electrónicos determinan los factores de
la ecuación \eqref{eq:el-formula-completa}. Esta composición conserva
la identidad de la sección a lo largo de sus distintas
representaciones, hasta la evaluación dimensional y el contraste
metrológico posterior.
